% =====================================================================
%  Origamizer: A Practical Algorithm for Folding Any Polyhedron
%  Algorithm diagrams (formal pseudocode) for every stage of the method.
%
%  Source paper: Erik D. Demaine and Tomohiro Tachi,
%    "Origamizer: A Practical Algorithm for Folding Any Polyhedron",
%    33rd Int. Symp. on Computational Geometry (SoCG 2017),
%    LIPIcs vol. 77, Article 34, pp. 34:1-34:16.
%    DOI: 10.4230/LIPIcs.SoCG.2017.34
%
%  Pseudocode authored from the prose description in the paper (the
%  paper itself gives no numbered pseudocode); section tags below point
%  back to the corresponding section of the original.
% =====================================================================
\documentclass[11pt]{article}

\usepackage[margin=1in]{geometry}
\usepackage{amsmath,amssymb}
\usepackage{xcolor}
\usepackage{enumitem}
\usepackage[ruled,vlined,linesnumbered]{algorithm2e}
\usepackage[hidelinks]{hyperref}

% ---- algorithm2e keyword setup -------------------------------------
\SetKwInput{KwData}{Input}
\SetKwInput{KwResult}{Output}
\SetKwComment{Comment}{$\triangleright$\ }{}
\SetKwFor{ForEach}{for each}{do}{end}
\SetKw{KwRet}{return}
\SetKwFunction{Origamizer}{Origamizer}
\SetKwFunction{TuckProxy}{ComputeTuckProxy}
\SetKwFunction{Squash}{SquashWafflePockets}
\SetKwFunction{Streams}{PlaceStreams}
\SetKwFunction{Sites}{PlaceSites}
\SetKwFunction{Fold}{VoronoiFolding}
\DontPrintSemicolon

% ---- shorthand -----------------------------------------------------
\newcommand{\eps}{\varepsilon}
\newcommand{\epsp}{\varepsilon'}
\newcommand{\dg}{^{\circ}}

\title{\textbf{The Origamizer Algorithm}\\[2pt]
\large Algorithm Diagrams for Folding Any Polyhedron}
\author{Pseudocode formalization of\\
Erik D.\ Demaine \& Tomohiro Tachi, \emph{Origamizer}, SoCG 2017}
\date{}

\begin{document}
\maketitle

\noindent
This document renders each algorithmic stage of the \emph{Origamizer}
method as explicit pseudocode. The original paper describes the
algorithm in prose (\S3--\S7); the section tags below cross-reference
those sections. Symbols follow the paper: $Q$ is the target oriented
\emph{polyhedral manifold} (disk topology, strictly convex facets, with
a \emph{clean} side and a \emph{tuck} side); $\eps>0$ is the requested
tolerance; $T$ is the \emph{tuck proxy} (a polyhedral \emph{waffle}
within $\eps$ of $Q$); and $\epsp$ with $0<\epsp<\eps$ lower-bounds how
far $T$ may extend beyond $Q$. The goal is a convex piece of paper $P$
and a \emph{watertight}, \emph{seamless}, \emph{$\eps$-extra} folding
$f\colon P\to Q$ with $f(P)\subseteq T$.

\vspace{1em}

% =====================================================================
\section*{Top-level driver}
% =====================================================================
\begin{algorithm}[H]
\caption[\textsc{Origamizer}]{\textsc{Origamizer}$(Q,\eps)$ --- fold a convex paper polygon into $Q$}
\KwData{Oriented polyhedral manifold $Q$ (disk, strictly convex facets,
        clean/tuck sides); tolerance $\eps>0$.}
\KwResult{Convex polygon $P$ and watertight, seamless, $\eps$-extra
          folding $f\colon P\to Q$ with $f(P)\subseteq T$.}
\BlankLine
$(T,\epsp)\leftarrow$ \TuckProxy{$Q,\eps$}\;
  \Comment*[r]{\S3: attach $\eps$-thin walls; split reflex curvature}
$\Sigma \leftarrow$ \Squash{$T,\epsp$}\;
  \Comment*[r]{\S4: flatten each waffle pocket; emit cells \& perpendiculars}
$(P',\,\mathcal{S},\,\delta)\leftarrow$ \Streams{$T,\epsp,\Sigma$}\;
  \Comment*[r]{\S5: Tutte-embed waffle dual; draw rivers \& brooks}
$(P,\,\mathcal{V})\leftarrow$ \Sites{$T,\epsp,\Sigma,\mathcal{S}$}\;
  \Comment*[r]{\S6: place point/segment/polygon sites; define $P$}
$f \leftarrow$ \Fold{$T,P,\mathcal{V}$}\;
  \Comment*[r]{\S7: fold $P$ into the tuck proxy via an abstract waffle}
\KwRet $(P,f)$\;
\end{algorithm}

% =====================================================================
\section*{\S3\quad Tuck Proxy and Waffles}
% =====================================================================
Computes the \emph{tuck proxy} $T$: a polyhedral waffle that contains
(and lies within $\eps$ of) $Q$. A waffle has \emph{floor} polygons
(forming a topological disk), \emph{wall quadrangles} glued along floor
edges, and \emph{wall triangles} glued at floor vertices. Insetting the
edges of $Q$ on the tuck side bisects every dihedral angle, splitting
each reflex (negative-curvature) angle into two convex ones, which makes
every \emph{waffle pocket} intrinsically convex.

\begin{algorithm}[H]
\caption[\textsc{ComputeTuckProxy}]{\textsc{ComputeTuckProxy}$(Q,\eps)$ --- build the tuck proxy waffle (\S3)}
\KwData{Polyhedral manifold $Q$; tolerance $\eps>0$.}
\KwResult{Tuck proxy $T$ with walls of height $\le\eps$; parameter
          $0<\epsp<\eps$ with no global self-intersection up to $\epsp$.}
\BlankLine
Floor polygons of $T \leftarrow$ the facets of $Q$\;
Choose $\epsp$ small enough that no collision event occurs while insetting\;
\ForEach{edge $e$ of $Q$ on the tuck side}{
  Inset $e$ inward by $\epsp$ (start of the 3D straight skeleton)\;
  \Comment*[r]{bisects the dihedral angle at $e$; a reflex angle splits into two convex angles}
}
\ForEach{vertex $v$ of $Q$}{
  On a small sphere around $v$, connect the incident edge offsets into a
  cycle\;
  \Comment*[r]{pocket bounded by floor polygon, two offset edges, one cycle edge; perimeter $\le 360\dg$}
  Overlay a regular tetrahedron on the cycle and triangulate if needed\;
  \Comment*[r]{makes the pocket intrinsically convex: edges $\le 109.5\dg$, perimeter $\le 328.5\dg<360\dg$}
}
Form wall quadrangles along floor edges and wall triangles at floor
vertices (height $\le\eps$)\;
\KwRet $(T,\epsp)$\;
\end{algorithm}

% =====================================================================
\section*{\S4\quad Waffle Pocket Squashing}
% =====================================================================
Locally ``squashes'' each waffle pocket flat into the plane by adding
material between polygons, only \emph{increasing} bottom angles and
preserving convexity, connectivity, and the cyclic order of wall
polygons. The flattened top edges form a convex \emph{cell}; its outward
edge normals are the \emph{cell perpendiculars} that will seed streams.

\begin{algorithm}[H]
\caption[\textsc{SquashWafflePockets}]{\textsc{SquashWafflePockets}$(T,\epsp)$ --- flatten every waffle pocket (\S4)}
\KwData{Tuck proxy $T$; parameter $\epsp$.}
\KwResult{For each pocket: a planar squashing (cell $+$ squashed wall
          polygons) and its cell perpendiculars.}
\BlankLine
\ForEach{waffle pocket $\pi$ (a floor vertex or a floor polygon)}{
  Vertex-unfold the wall polygons of $\pi$ into the plane, spreading the
  angular gaps evenly between squashed wall polygons\;
  \ForEach{gap}{
    Place a ray from the vertex along the angular bisector of the two
    adjacent squashed wall polygons (rounded to a polygon boundary if
    needed)\;
    \Comment*[r]{consecutive rays sandwich each wall polygon at an angle $<180\dg$}
  }
  \uIf{$\pi$ is a floor vertex}{
    Cell $\leftarrow$ a small \emph{circle} surrounding the vertex\;
  }
  \Else(\tcp*[h]{$\pi$ is a floor polygon}){
    Cell $\leftarrow$ offset of the floor polygon, corners smoothed with
    quadratic B\'ezier arcs\;
    \Comment*[r]{wall quadrangles then squash into trapezoids}
  }
  Clip the rays against the cell to bound the squashed wall polygons
  (strictly convex; bottom angles only increase)\;
  \ForEach{cell edge}{
    Emit a cell perpendicular: outward normal starting at the midpoint of
    the bottom of the corresponding squashed wall polygon, oriented CCW\;
  }
}
\KwRet the collection of squashings and cell perpendiculars\;
\end{algorithm}

% =====================================================================
\section*{\S5\quad Placing Streams}
% =====================================================================
Embeds the \emph{waffle dual} in the plane (Tutte embedding) inside a
convex polygon $P'$, then connects the placed floor vertices/polygons by
\emph{streams}: thick \emph{rivers} for wall quadrangles and zero-width
\emph{brooks} for wall triangles. Three nested disks around each node
re-aim the stream directions to match the embedded dual without
collisions; filleting makes streams $C^1$ with positive clearance
$\delta$.

\begin{algorithm}[H]
\caption[\textsc{PlaceStreams}]{\textsc{PlaceStreams}$(T,\epsp,\Sigma)$ --- route rivers and brooks (\S5)}
\KwData{Tuck proxy $T$; parameter $\epsp$; squashing $\Sigma$.}
\KwResult{Convex polygon $P'$; stream placement $\mathcal{S}$; clearance
          $\delta>0$.}
\BlankLine
\textbf{Part 1 --- embed the waffle dual.}\;
Compute a Tutte embedding of the waffle dual with its outside face, so
boundary nodes lie on $\partial P'$ ($P'$ convex)\;
Ensure every river-attached boundary node has room to be thickened
parallel to its edge of $P'$ without leaving it or hitting another
boundary node\;
\BlankLine
\textbf{Part 2 --- place nodes and connect.}\;
\ForEach{waffle dual node $x$}{
  \lIf{$x$ is a floor vertex}{place a point}
  \lElse{place an isometric copy of the corresponding facet of $Q$}
}
\ForEach{waffle dual edge}{
  \uIf{it is a wall quadrangle}{
    Draw a \emph{river}: connect two equal-length floor-polygon edges and
    thicken by that edge length\;
  }\Else{
    Draw a \emph{brook}: connect two points with a zero-width $C^1$ curve\;
  }
}
\ForEach{node $x$ of degree $k$}{
  Build nested disks $R_x\subset S_x\subset T_x$ (radii $r_x<s_x<t_x$):\;
  \quad $R_x$ separates streams from the floor vertex/polygon\;
  \quad $S_x$ bends streams to meet $\partial S_x$ orthogonally\;
  \quad $T_x$ twists streams onto incident dual-edge directions using $k$
        collision-free tracks\;
  Rotate the floor polygon so one edge normal needs no twist; cut the
  $T_x\!\setminus\!S_x$ annulus there\;
  Route each connection inside$\to$outside in CCW order: outermost free
  track for leftward (CCW) links, innermost free track for rightward (CW)
  links\;
}
Fillet every sharp corner with a circular arc $\Rightarrow$ streams are
$C^1$ with positive clearance $\delta$\;
Scale up the Tutte embedding by a finite factor so all disks stay local\;
\KwRet $(P',\mathcal{S},\delta)$\;
\end{algorithm}

% =====================================================================
\section*{\S6\quad Placing Sites}
% =====================================================================
Places point, segment, and polygon \emph{sites} on the paper so that
their generalized Voronoi diagram --- taken in the \emph{geodesic vertex
metric} --- contracts to the waffle graph of $T$. Protective regions
called \emph{squashes} (pumpkins, jack-o'-lanterns, calabashes,
cucumbers) keep new sites from corrupting earlier structure. The piece
of paper $P$ is then the convex hull of the resulting paired subcells.

\begin{algorithm}[H]
\caption[\textsc{PlaceSites}]{\textsc{PlaceSites}$(T,\epsp,\Sigma,\mathcal{S})$ --- site placement (\S6)}
\KwData{Tuck proxy $T$; parameter $\epsp$; squashing $\Sigma$; streams
        $\mathcal{S}$.}
\KwResult{Piece of paper $P$ and a set $\mathcal{V}$ of point/segment/%
          polygon sites.}
\BlankLine
\textbf{(1) Node sites.}\;
\ForEach{brook node}{place a \emph{pumpkin} (point site cluster)}
\ForEach{placed floor polygon}{place a \emph{jack-o'-lantern}\;
  \Comment*[r]{core $=$ cell of the local squashing of that pocket}}
\BlankLine
\textbf{(2) Stream sites.}\;
\ForEach{brook}{
  Lay point sites along it; insert a \emph{calabash} between consecutive
  sites with flesh-polygon angle $=$ bottom angle of the brook's wall
  triangle\;
}
\ForEach{river}{
  Lay segment sites along it; insert a \emph{cucumber} between consecutive
  sites with flesh-polygon angles $=$ bottom angles of the river's wall
  quadrangle\;
}
\If{a gap between consecutive sites lets another stream hit the squash}{
  Bisect the gap and add a site; subdivide until every squash point is
  $\le\tfrac12\epsp$ from the core\;
}
\BlankLine
\textbf{(3) Bank sites.}\;
\While{some paper point is $>\epsp$ (Euclidean) from its nearest site}{
  Add a point site there\;
  \Comment*[r]{squashes $\subseteq$ Minkowski sum of each site with an $\epsp$-disk $\Rightarrow$ terminates off the squashes}
}
\BlankLine
\textbf{(4) Define the paper.}\;
$P\leftarrow$ convex hull (within $P'$) of the paired subcells\;
\Comment*[r]{equivalently: convex hull of the sites and the Voronoi diagram, clipped to $P'$}
\KwRet $(P,\mathcal{V})$\;
\end{algorithm}

\noindent\emph{Invariants guaranteed by construction.}
(i) the Voronoi diagram contracts onto the modified dual --- the waffle
graph of $T$; (ii) each Voronoi edge is straight and mirror-reflects its
defining site vertices, forming \emph{paired subcells} (quadrangular iff
the edge realizes a wall quadrangle); (iii) every realizing Voronoi
edge's subcell angles are $\ge$ the bottom angles of the corresponding
wall polygon; (iv) every point of $P$ lies within $\epsp$ of a site
vertex.

% =====================================================================
\section*{\S7\quad Voronoi Folding}
% =====================================================================
Folds $P$ into the tuck proxy $T$ by a sequence of \emph{folding steps},
each producing an \emph{abstract waffle} (a metric polyhedral complex)
and never separating layers already brought together. The final abstract
waffle embeds isometrically onto $T$, giving the desired folding $f$.

\begin{algorithm}[H]
\caption[\textsc{VoronoiFolding}]{\textsc{VoronoiFolding}$(T,P,\mathcal{V})$ --- fold $P$ onto the tuck proxy (\S7)}
\KwData{Tuck proxy $T$; paper $P$; sites/Voronoi $\mathcal{V}$.}
\KwResult{Watertight, seamless, $\eps$-extra folding $f\colon P\to Q$
          with $f(P)\subseteq T$.}
\BlankLine
\textbf{(1) Initial fold into an abstract waffle $A$.}\;
$G\leftarrow\{$Voronoi edge of each paired subcell$\}\cup\{$boundary edge
of each unpaired subcell$\}$\;
\ForEach{Voronoi edge of $G$}{
  Fold its two (mirror-image) paired subcells onto each other to
  \emph{doubly} cover the matching wall polygon of $A$\;
}
\ForEach{boundary edge of $G$}{
  Fold the unpaired subcell to \emph{singly} cover its wall polygon\;
}
\Comment*[r]{floor of $A$ $=$ placed floor polygons $\Rightarrow$ isometric to $Q$; each Voronoi cell $\to$ a pocket of $A$}
\BlankLine
\textbf{(2) Contract \& merge non-realizing edges.}\;
\ForEach{Voronoi edge not realizing a waffle-graph edge of $T$}{
  Contract it: fold the wall triangle in half along its bisector and glue
  it onto an adjacent wall polygon of the same pocket\;
  \If{several edges join the same two nodes}{
    Crimp-fold the larger angles down to the smallest, stack the now
    isometric wall polygons, and collapse to a pocket\;
  }
}
\Comment*[r]{leaves a simplified waffle graph whose dual is just the realizing paths}
\BlankLine
\textbf{(3) Merge realizing edges.}\;
\ForEach{dual path (sequence of wall polygons for one waffle-graph edge)}{
  Crimp-fold their bottom angles to match the bottom angles of the
  corresponding wall polygon of $T$\;
  Stack them as a single wall polygon and embed it isometrically into
  that wall polygon of $T$\;
}
$f\leftarrow$ the resulting folded state $\subseteq T$\;
\KwRet $f$\;
\end{algorithm}

\noindent\emph{Correctness of the output.}
\textbf{Seamless:} the floor polygons are isometric to $Q$ and no other
paper folds strictly interior to a facet. \textbf{Watertight:} as a point
$p$ moves CCW along $\partial P$ and $q$ moves CCW along $\partial Q$,
the correspondence $f(p)=q$ on boundary sites (and $q$ pinned to the
nearest boundary vertex between sites) stays within Fr\'echet distance
$\eps$. \textbf{$\eps$-extra:} every added wall feature has height $\le\eps$
and $f(P)\subseteq T$, so the Hausdorff distance from $Q$ is $\le\eps$.

\end{document}
